\ifintegriereaufgaben
\section*{Potenzreihen Aufgaben}\label{MA2-potenzreihen-aufgaben}

Erkl�rungen und Beispiele, die zur Bearbeitung der Aufgaben hilfreich sind, finden Sie im Skript\wennreferenzvorhanden{MA2-potenzreihen}{ (Abschnitt \abschnittsverweis{MA2-potenzreihen})} und, falls Sie nicht weiterkommen, dann schauen Sie \hyperref[MA2-potenzreihen-loesungen]{hier}.
\else
\section{Potenzreihen Aufgaben}\label{MA2-potenzreihen-aufgaben}

Regeln und Beispiele zu diesen Aufgaben finden Sie im Skript. Die L�sungen finden Sie \hyperref[MA2-potenzreihen-loesungen]{hier}.
\fi

\begin{exercisebox}[Konvergenz] %Papula VI, Abschnitt 2, Aufgabe 1
	Bestimmen Sie den Konvergenzradius und geben Sie den Konvergenzbereich der folgenden Reihen an. Untersuchen Sie dabei auch die R�nder des Konvergenzbereichs.
	\begin{eqnarray*}
		\begin{array}{l l l}
			(a)\; P(x) = \displaystyle{\sum\limits_{k=1}^\infty} k\cdot x^k \quad & (b)\; P(x) = \displaystyle{\sum\limits_{k=1}^\infty} (-1)^k \cdot \displaystyle{\frac{x^k}{k}} \quad & (c)\; P(x) = \displaystyle{\sum\limits_{k=1}^\infty} \displaystyle{\frac{x^k}{k^2}} \\
			(d)\; P(x) = \displaystyle{\sum\limits_{k=0}^\infty} \displaystyle{\frac{x^k}{2^k}} \quad & (e)\; P(x) = \displaystyle{\sum\limits_{k=0}^\infty } \displaystyle{\frac{k}{k+1} x^{k+1}} \quad & (f)\; P(x) = \displaystyle{\sum\limits_{k=0}^\infty} \displaystyle{\frac{k+1}{k!}} x^k
		\end{array}
	\end{eqnarray*}
\end{exercisebox}

%Moodle %HMI I, Blatt 11, Aufgabe 1 Moodle
\begin{exercisebox}[Konvergenz] %HMI I, Blatt 10, Aufgabe 4 Moodle
	Bestimmen Sie den Konvergenzradius der folgenden Reihen.
	\begin{eqnarray*}
		(a)\; \displaystyle{\sum\limits_{k=1}^\infty} \left(  \displaystyle{\frac{k!}{3 \cdot 5\cdot \dots \cdot (2k+1)}}\right)^{2} x^k \quad  (b)\; \displaystyle{\sum\limits_{k=1}^\infty} \left( \begin{array}{c} 2k\\ k\end{array}\right) x^k \quad  (c)\; \displaystyle{\sum\limits_{k=1}^\infty } \left( k^4 - 4k^3\right) x^k
	\end{eqnarray*}
\end{exercisebox}

\begin{exercisebox}[ Analytische Funktionen]
Begr�nden Sie, warum die folgende Funktion nicht analytisch ist
\[ f: x\mapsto  \begin{cases} e^{-1/x^2}, \quad x\not=0\,,\\
0, \quad x=0\,.\end{cases}
\]
\end{exercisebox}

\begin{exercisebox}[Leibnitz-Kriterium] %HMI I, Blatt 10, Aufgabe 2
	Der Weg eines Teilchens in Abh�ngigkeit der Beschleunigung $x<0$ sei durch die unendliche Reihe
	\[
		S(x) = \sum\limits_{k=1}^\infty \frac{x^{k-1}}{k^2\cdot 3^k}
	\] gegeben.  Untersuchen Sie mit Hilfe des Leibnizkriteriums f�r welche $x\in\mathbb R$ die Reihe $S(x)$ sinnvoll gegeben ist.
\end{exercisebox}

\begin{exercisebox}[Potenzreihen und Differentialgleichungen] %Alternativ: Blatt 11, Aufgabe 3
  Ein K�rper mit der Masse $m$ bewege sich auf einer geradlinigen Bahn. Die Anfangsgeschwindigkeit sei $v_0 = 1$. Durch eine Kraft $F$, die der Geschwindigkeit proportional ist, wird der K�rper abgebremst.

  \begin{itemize}
	  \item[(a)] Stellen Sie  Bewegungsgleichung f�r den K�rper auf.
	  \item[(b)] L�sen Sie die Differentialgleichung mit Hilfe eines Potenzreihenansatzes.
  \end{itemize}
\end{exercisebox}

%
%\begin{exercisebox}[Ausgleichsgerade]
%  Eine Person ist $1,80 \mathrm{m}$ gro�. Sie entfernt sich auf einer waagerecht verlaufenden Stra�e mit der Geschwindigkeit $6\frac{\mathrm{km}}{\mathrm{h}}$ von einer Lichtquelle. Die Lichtquelle befindet sich in einer H�he von $3,50\mathrm{m}$. Mit welcher Geschwindigkeit wandert der Schatten des Kopfes auf der Stra�e entlang?
%\end{exercisebox}

%\begin{mybox}(Ideen zur Approximation von Integralen)
%	Polynome zu gewissem Grad exakt integrieren.
%\end{mybox}
%
%\begin{mybox}(Lineare Differentialgleichungen: Approximationen von h�herer Ordnung)
%\end{mybox}
%
% Minimierungsproblem in $\mathbb R^n$
%		\[
%		x\in \mathbb R\setminus\{ 0\}:\quad \frac{\mathrm{d}}{\mathrm{d} x }  |x|= \frac{\mathrm{d} }{\mathrm{d} x } \sqrt{x^2} = \frac12\cdot (x^2)^{-\frac12} \cdot 2x = \frac{x}{\sqrt{x^2}} = \frac{x}{|x|} = \begin{cases} 1, \quad x > 0\,,  1, \quad x<0 \end{cases} \,.
%	\]


